\chapter{Cel i zakres pracy}
Jasne wyznaczenie kierunków działania pozwala na precyzyjne zaplanowanie etapów pracy i uniknięcie zbędnych funkcjonalności. Dokumentacja ta definiuje granice projektu oraz to, co dokładnie zostanie dostarczone po jego zakończeniu.
\section{Cel projektu}
Głównym celem projektu było opracowanie i wdrożenie kompletnego, rozproszonego systemu pozwalającego na zdalny odczyt, archiwizację i wizualizację danych pomiarowych z fizycznego czujnika temperatury. System miał łączyć w sobie elementy sprzętowe, serwerowe oraz klienckie, tworząc spójną całość gotową do działania w warunkach lokalnych.

\section{Cele szczegółowe}
W~ramach realizacji projektu wyznaczono następujące cele szczegółowe:
\begin{itemize}
    \item zaprogramowanie mikrokontrolera ESP32 do obsługi czujnika pomiarowego oraz bezprzewodowego przesyłania danych przez Wi-Fi,
    \item zaimplementowanie serwera backendowego w technologii Spring Boot, pełniącego rolę centralnego węzła komunikacyjnego,
    \item konfiguracja i konteneryzacja relacyjnej bazy danych MySQL do trwałego przechowywania historii pomiarów,
    \item stworzenie natywnej aplikacji mobilnej na system Android, umożliwiającej bieżący podgląd danych,
    \item budowa responsywnego interfejsu webowego w formie strony internetowej do przeglądania dziennika zapisów,
    \item integracja wszystkich modułów za pomocą protokołu HTTP i formatu danych JSON.
\end{itemize}

\section{Zakres pracy}
Zakres pracy obejmuje następujące kroki:
\begin{itemize}
    \item konfiguracja warstwy sprzętowej obejmującej mikrokontroler ESP32 oraz czujnik temperatury,
    \item implementacja logiki serwerowej wraz z obsługą endpointów REST API,
    \item przygotowanie schematu bazy danych oraz skryptów łączących serwer z bazą MySQL,
    \item opracowanie interfejsu graficznego i logiki komunikacyjnej dla aplikacji mobilnej Android,
    \item przygotowanie frontendu webowego (HTML/CSS/JS) współpracującego z API serwera,
    \item przeprowadzenie testów funkcjonalnych i weryfikacja poprawności przepływu danych w sieci lokalnej.
\end{itemize}